\section{收敛性结果与完整证明}\label{sec:theory}
\subsection{假设、随机信息与评价指标}
以下只分析算法\ref{alg:exact}。令 $H_t=\nabla f(W_t)$，$Z_t=G_t-H_t$，$E_t=M_t-H_t$。用 $\mathcal F_t$ 表示采样 $G_t$ 之前已知的全部信息，其中包括 $W_t$ 和过去的样本。

\begin{assumption}[下界与光滑性]\label{ass:smooth}
存在有限 $f_{\inf}$，使 $f(W)\geq f_{\inf}$。对任意 $X,Y$，有
\begin{equation}\label{eq:smooth}
 \Frob{\nabla f(X)-\nabla f(Y)}\leq L\Frob{X-Y},\qquad L>0.
\end{equation}
记 $\Delta=f(W_1)-f_{\inf}$，初值 $W_1$ 固定。
\end{assumption}

\begin{assumption}[条件无偏与方差]\label{ass:noise}
对每个 $t$，批量梯度满足
\begin{equation}\label{eq:noise}
 \E[Z_t\mid\mathcal F_t]=0,\qquad
 \E[\Frob{Z_t}^2\mid\mathcal F_t]\leq\sigma^2/B.
\end{equation}
\end{assumption}

若每个样本的条件方差不超过 $\sigma^2$ 且同一批中的样本条件独立，则平均梯度满足式\eqref{eq:noise}。证明仅需要这一矩条件，不要求噪声服从高斯分布，也不假定不同时间的参数或真实梯度相互独立。合成实验采用高斯噪声只是为了精确控制其方差。

本文评价平均梯度核范数
\begin{equation}\label{eq:criterion}
 A_T=\frac1T\sum_{t=1}^T\E\nuc{H_t}.
\end{equation}
由于 $\Frob{H_t}\leq\nuc{H_t}$，$A_T$ 的上界也控制平均 Frobenius 梯度范数。若独立均匀抽取 $R\in\{1,\ldots,T\}$，则 $\E\Frob{\nabla f(W_R)}\leq A_T$。这给出随机输出的一阶平稳性保证，不要求最后一个迭代点具有同样上界。

\subsection{含跟踪误差的下降不等式}
\begin{lemma}[单步下降]\label{lem:descent}
在假设\ref{ass:smooth}下，算法\ref{alg:exact}的每一步都满足
\begin{equation}\label{eq:descent}
 f(W_{t+1})\leq f(W_t)-\eta\nuc{H_t}
 +2\eta\sqrt r\Frob{E_t}+\frac{Lr\eta^2}{2}.
\end{equation}
因此
\begin{equation}\label{eq:master}
 A_T\leq\frac{\Delta}{\eta T}
 +\frac{2\sqrt r}{T}\sum_{t=1}^T\E\Frob{E_t}
 +\frac{Lr\eta}{2}.
\end{equation}
\end{lemma}
\begin{proof}
由谱范数与核范数的 H\"older 不等式，以及引理\ref{lem:duality}，
\begin{align}
 \inner{H_t}{O_t^\star}
 &=\inner{M_t}{O_t^\star}-\inner{E_t}{O_t^\star} \notag\\
 &\geq\nuc{M_t}-\nuc{E_t}\op{O_t^\star}\notag\\
 &\geq\nuc{H_t}-2\nuc{E_t}\notag\\
 &\geq\nuc{H_t}-2\sqrt r\Frob{E_t}.\label{eq:alignment}
\end{align}
第三行使用核范数的三角不等式；这里没有使用 $M_t$ 的逆或其最小奇异值，因此秩亏情形仍成立。光滑性推出标准下降引理
\begin{equation}
 f(W_{t+1})\leq f(W_t)+\inner{H_t}{W_{t+1}-W_t}
 +\frac L2\Frob{W_{t+1}-W_t}^2.
\end{equation}
代入 $W_{t+1}-W_t=-\eta O_t^\star$，并使用 $\Frob{O_t^\star}^2\leq r$，得到式\eqref{eq:descent}。对 $t$ 求和、取期望，再用 $f(W_{T+1})\geq f_{\inf}$，即得式\eqref{eq:master}。
\end{proof}

式\eqref{eq:descent}中的三个变化项具有不同含义：核范数项产生下降，跟踪误差造成方向偏差，步长平方项来自目标曲率。仅有动量矩阵的范数较小并不足以保证下降；真正相关的是它与当前梯度的差。

\subsection{动量跟踪误差的三项分解}
\begin{lemma}[误差递推与平均界]\label{lem:tracking}
在两个假设下，令 $\alpha=1-\beta\in(0,1]$，则对 $t\geq1$，
\begin{equation}\label{eq:trackingpoint}
 \E\Frob{E_t}\leq
 \frac{\sigma}{\sqrt B}\beta^{t-1}
 +\frac{\sigma}{\sqrt B}\sqrt{\frac{\alpha}{1+\beta}}
 +\frac{\beta L\eta\sqrt r}{\alpha}.
\end{equation}
其中在 $t=1$ 时将 $\beta^0$ 解释为 1。进一步有
\begin{equation}\label{eq:trackingavg}
 \frac1T\sum_{t=1}^T\E\Frob{E_t}\leq
 \frac{\sigma}{T\alpha\sqrt B}
 +\sigma\sqrt{\frac{\alpha}{B}}
 +\frac{\beta L\eta\sqrt r}{\alpha}.
\end{equation}
\end{lemma}
\begin{proof}
初始化给出 $E_1=Z_1$。对 $t\geq2$，由动量递推可得
\begin{equation}\label{eq:erecursion}
 E_t=\beta E_{t-1}+\alpha Z_t+\beta(H_{t-1}-H_t).
\end{equation}
将递推展开为
\begin{equation}\label{eq:edecomposition}
 E_t=\underbrace{\beta^{t-1}Z_1}_{\text{初始化}}
 +\underbrace{\alpha\sum_{i=2}^t\beta^{t-i}Z_i}_{\text{噪声}}
 +\underbrace{\beta\sum_{i=2}^t\beta^{t-i}(H_{i-1}-H_i)}_{\text{梯度漂移}}.
\end{equation}
先由 Jensen 不等式，初始化项的期望范数不超过 $\beta^{t-1}\sigma/\sqrt B$。

对于噪声项，当 $i<j$ 时，$Z_i$ 对 $\mathcal F_j$ 可测，故
\begin{equation}
 \E\inner{Z_i}{Z_j}
 =\E\inner{Z_i}{\E[Z_j\mid\mathcal F_j]}=0.
\end{equation}
这一步利用的是鞅差正交性，而非未经说明的时间独立性。于是
\begin{align}
 \E\Frob{\alpha\sum_{i=2}^t\beta^{t-i}Z_i}
 &\leq\left(\alpha^2\sum_{i=2}^t\beta^{2(t-i)}\E\Frob{Z_i}^2\right)^{1/2}\notag\\
 &\leq\frac{\alpha\sigma}{\sqrt{B(1-\beta^2)}}
 =\frac{\sigma}{\sqrt B}\sqrt{\frac{\alpha}{1+\beta}}.\label{eq:noisesum}
\end{align}
当 $t=1$ 时，该和为空，右端仍是有效上界。

最后，由光滑性与 $\Frob{O_{i-1}^\star}\leq\sqrt r$，有
\begin{equation}
 \Frob{H_{i-1}-H_i}\leq L\Frob{W_{i-1}-W_i}
 \leq L\eta\sqrt r.
\end{equation}
漂移项的范数因此不超过
\begin{equation}\label{eq:drift}
 \beta L\eta\sqrt r\sum_{i=2}^t\beta^{t-i}
 \leq\frac{\beta L\eta\sqrt r}{\alpha}.
\end{equation}
将三项上界代入式\eqref{eq:edecomposition}，得到式\eqref{eq:trackingpoint}。再用 $\sum_{t=1}^T\beta^{t-1}\leq1/\alpha$ 和 $1+\beta\geq1$，得到式\eqref{eq:trackingavg}。
\end{proof}

减小 $\alpha$ 会增强动量平均，从而减小噪声项，但也使漂移项中的 $1/\alpha$ 增大。因此，动量参数并非越接近 1 越好；必须同时缩小参数移动的步长。这个权衡直接决定下面的随机参数调度。

\subsection{有限步界与确定性收敛率}
\begin{theorem}[显式有限步上界]\label{thm:finite}
在假设\ref{ass:smooth}与\ref{ass:noise}下，对任意 $T\geq1$、$\eta>0$、$0\leq\beta<1$，算法\ref{alg:exact}满足
\begin{equation}\label{eq:finite}
 \boxed{\begin{aligned}
 A_T\leq {}&\frac{\Delta}{\eta T}+\frac{Lr\eta}{2}\\
 &+2\sqrt r\left(
 \frac{\sigma}{T\alpha\sqrt B}
 +\sigma\sqrt{\frac{\alpha}{B}}
 +\frac{\beta L\eta\sqrt r}{\alpha}\right).
 \end{aligned}}
\end{equation}
\end{theorem}
\begin{proof}
将引理\ref{lem:tracking}的平均误差界代入引理\ref{lem:descent}的式\eqref{eq:master}即可。由于所有被估计量均为非负且有限的范数矩，本证明中的有限求和与期望交换是合法的。
\end{proof}

\begin{corollary}[确定性情形]\label{cor:det}
若 $G_t=H_t$，固定 $\beta<1$，令
\begin{equation}
 C_\beta=Lr\left(\frac12+\frac{2\beta}{1-\beta}\right),
\end{equation}
则 $T^{-1}\sum_t\nuc{H_t}\leq\Delta/(\eta T)+C_\beta\eta$。
当 $\Delta>0$ 且 $\eta=\sqrt{\Delta/(C_\beta T)}$ 时，
\begin{equation}\label{eq:detrate}
 \frac1T\sum_{t=1}^T\Frob{H_t}
 \leq\frac1T\sum_{t=1}^T\nuc{H_t}
 \leq2\sqrt{\frac{\Delta C_\beta}{T}}.
\end{equation}
\end{corollary}
\begin{proof}
取 $\sigma=0$，将式\eqref{eq:finite}中所有与 $\eta$ 成正比的项合并即得 $C_\beta\eta$。对 $\Delta/(\eta T)+C_\beta\eta$ 最小化得到所给步长。若 $\Delta=0$，初值达到目标下界，光滑无约束目标的梯度为零，精确算法保持不动。
\end{proof}

特别地，$\beta=0$ 时没有跟踪误差，$C_0=Lr/2$。固定非零步长只能由此保证一个随步长缩小的平均残差上界，不能据此宣称参数序列必然收敛。步长按 $T^{-1/2}$ 缩放指的是针对不同预算分别设置常步长，并非本定理已经分析了同一次运行中的逐步衰减。

\subsection{随机非凸收敛率与梯度复杂度}
\begin{theorem}[固定批量的随机收敛率]\label{thm:stochastic}
固定 $B\geq1$。对给定迭代预算 $T$，取
\begin{equation}\label{eq:schedule}
 \alpha=T^{-1/2},\qquad \beta=1-T^{-1/2},\qquad
 \eta=cT^{-3/4},\quad c>0.
\end{equation}
在上述假设下，有
\begin{equation}\label{eq:stochrate}
 \begin{aligned}
 A_T\leq{}&\left(\frac\Delta c+2Lrc+\frac{2\sqrt r\sigma}{\sqrt B}\right)T^{-1/4}\\
 &+\frac{2\sqrt r\sigma}{\sqrt B}T^{-1/2}
 +\frac{Lrc}{2}T^{-3/4}.
 \end{aligned}
\end{equation}
因此，在 $L,r,\Delta,\sigma,B,c$ 固定时，$A_T=O(T^{-1/4})$。
\end{theorem}
\begin{proof}
将式\eqref{eq:schedule}逐项代入式\eqref{eq:finite}。第一项变为 $(\Delta/c)T^{-1/4}$，曲率项为 $(Lrc/2)T^{-3/4}$，初始化项为 $(2\sqrt r\sigma/\sqrt B)T^{-1/2}$，噪声项为 $(2\sqrt r\sigma/\sqrt B)T^{-1/4}$。漂移项为 $2\beta LrcT^{-1/4}$，利用 $\beta\leq1$ 即得到式\eqref{eq:stochrate}。
\end{proof}

由式\eqref{eq:stochrate}，若要求均匀随机输出满足 $\E\Frob{\nabla f(W_R)}\leq\varepsilon$，足够的迭代量为 $O(\varepsilon^{-4})$，单样本梯度调用量为 $O(B\varepsilon^{-4})$。固定 $B$ 时二者具有相同的 $\varepsilon$ 阶数。这里的指标是梯度\emph{范数}的期望，而非梯度范数\emph{平方}的期望，两种复杂度表述不能直接互换。

该界保留了 $r$ 的显式依赖。它并未证明 Muon 的常数小于 AdamW，也未证明其为一般随机优化问题的最优方法。若在随机情形固定 $\beta<1$ 和 $B$，式\eqref{eq:finite}中的 $\sigma\sqrt{\alpha/B}$ 项不会随 $T$ 消失，当前证明只能给出噪声残差界。应用中常用的固定动量因此不属于定理\ref{thm:stochastic}的消失误差调度。

\subsection{含谱范数近似误差的完整有限步定理}\label{sec:inexact-theorem}
为使近似正交化进入主结论，以下将四项条件分别列明：目标有有限下界；满足式\eqref{eq:smooth}的 Frobenius 光滑性；随机梯度满足式\eqref{eq:noise}；实际方向与理想极因子的谱范数距离有统一上界。前三项沿用前文假设，第四项为
\begin{equation}\label{eq:delta-assumption}
 O_t^\star=\polar(M_t),\qquad
 \|O_t-O_t^\star\|_2\leq\delta\leq\bar\delta<1,
 \quad\text{对全部 }t\text{ 几乎处处成立}.
\end{equation}
这里 $O_t$ 是未作形状缩放的近似方向，更新为 $W_{t+1}=W_t-\eta O_t$；动量仍采用算法\ref{alg:exact}的初始化和普通 EMA。该条件定理不同时包含权重衰减、Nesterov 修正或辅助 AdamW。

\begin{theorem}[带谱范数近似误差的非渐近界]\label{thm:inexact}
在上述四项条件下，令 $\alpha=1-\beta$，则
\begin{equation}\label{eq:inexact-finite}
\boxed{\begin{aligned}
 A_T\leq\frac1{1-\delta}\biggl[&\frac\Delta{\eta T}
 +\frac{Lr\eta(1+\delta)^2}{2}\\
 &+2\sqrt r\left(\frac\sigma{T\alpha\sqrt B}
 +\sigma\sqrt{\frac\alpha B}
 +\frac{\beta L\eta\sqrt r(1+\delta)}\alpha\right)\biggr].
\end{aligned}}
\end{equation}
固定 $B$，取 $\eta=cT^{-3/4}$、$\alpha=T^{-1/2}$，且 $\delta\leq\bar\delta<1$ 与 $T$ 无关，则 $A_T=O(T^{-1/4})$。特别地，$\delta=O(T^{-1/4})$ 并最终小于 $\bar\delta$ 也是充分条件。
\end{theorem}
\begin{proof}
设 $R_t=O_t-O_t^\star$。由谱范数与核范数对偶关系，
\begin{align}
 \inner{H_t}{O_t}
 &=\inner{H_t}{O_t^\star}+\inner{H_t}{R_t}\notag\\
 &\geq(1-\delta)\nuc{H_t}-2\sqrt r\Frob{E_t}.
\end{align}
并且 $\|O_t\|_2\leq1+\delta$，所以 $\|O_t\|_F\leq\sqrt r(1+\delta)$。代入下降引理得
\begin{equation}\label{eq:inexact-descent}
 f(W_{t+1})\leq f(W_t)-\eta(1-\delta)\nuc{H_t}
 +2\eta\sqrt r\Frob{E_t}+\frac{Lr\eta^2(1+\delta)^2}{2}.
\end{equation}
式\eqref{eq:erecursion}的代数递推保持不变；噪声和初始化项也不变。实际更新的长度改变了漂移界：
\begin{equation}
 \|H_{t-1}-H_t\|_F\leq L\eta\sqrt r(1+\delta).
\end{equation}
因此式\eqref{eq:trackingavg}的最后一项变为 $\beta L\eta\sqrt r(1+\delta)/\alpha$。将这一平均误差界代入式\eqref{eq:inexact-descent}的望远镜求和，利用目标下界，再除以 $\eta T(1-\delta)$，即得式\eqref{eq:inexact-finite}。代入所给调度后，初始目标差项、噪声项和漂移项均不超过常数乘 $T^{-1/4}$，其余项阶数更高；$1/(1-\delta)\leq1/(1-\bar\delta)$，所以得到所述速率。
\end{proof}

\paragraph{为什么不必强制 $\delta\to0$？}
谱范数误差通过 $|\langle H_t,R_t\rangle|\leq\delta\|H_t\|_*$ 相对吸收进下降项。因此统一的 $\delta<1$ 已是这条证明路线的充分条件，而不是必须额外假定 $\delta=O(T^{-1/4})$。当 $\delta=0$ 时，式\eqref{eq:inexact-finite}退化为精确版本的界。若 $\delta\geq1$，该证明失去正下降系数；这只表示本充分界不可用，并不说明实际算法必然发散。

\subsection{谱范数光滑性、固定步数与适用范围}
若采用谱范数几何，梯度的自然对偶范数为核范数。一个明确的充分光滑条件是
\begin{equation}\label{eq:spectral-smooth}
 \|\nabla f(X)-\nabla f(Y)\|_*\leq L_s\|X-Y\|_2.
\end{equation}
它推出以 $L_s\|Y-X\|_2^2/2$ 为二次项的下降引理。重复上述证明，曲率项不再含 $r$，而梯度漂移用 $\|H_t-H_{t-1}\|_F\leq\|H_t-H_{t-1}\|_*\leq L_s\eta(1+\delta)$ 控制，得到明确的谱范数几何版本：
\begin{equation}\label{eq:spectral-inexact-finite}
\begin{aligned}
 A_T\leq\frac1{1-\delta}\biggl[&\frac\Delta{\eta T}+\frac{L_s\eta(1+\delta)^2}{2}\\
 &+2\sqrt r\left(\frac\sigma{T\alpha\sqrt B}+\sigma\sqrt{\frac\alpha B}
 +\frac{\beta L_s\eta(1+\delta)}\alpha\right)\biggr].
\end{aligned}
\end{equation}
外部的 $2\sqrt r$ 仍由方差所用的 Frobenius 范数转换而来。同样代入 $\eta=cT^{-3/4}$、$\alpha=T^{-1/2}$，得到平均核范数梯度的 $O(T^{-1/4})$ 界。不能只写“谱范数光滑”而不说明梯度在哪个对偶范数下计量。Frobenius 光滑可给出保守常数 $L_s=rL$，所以这两种写法不会自动带来维度上的改进。

Higham\citep{higham2008}讨论了极分解与矩阵函数迭代；本文五次多项式的具体系数取自 Muon 的工程实现\citep{jordan2024}，两者不应混作同一收敛保证。对固定 $K$，合成标量多项式 $\phi^{\circ K}(s)$ 在 $s\downarrow0$ 时趋于零，而理想极因子在非零支撑上的奇异值是 1。因此，在任意病态谱上不能仅由“执行五步”推出一个与输入无关的 $\bar\delta<1$。若要把 $\delta$ 的预算调度落实为 $K$ 的调度，还需要谱条件与多项式误差分析。

更直接地，经验五次多项式满足
\begin{equation}
 \phi(1)=3.4445-4.7750+2.0315=0.7010\ne1.
\end{equation}
所以 1 不是这个固定多项式的不动点。由连续性，反复使用同一组系数的奇异值序列不能以 1 为极限；单纯增大 $K$ 不会自动使谱范数误差趋于零。若需要可趋零的精确极因子误差保证，就必须采用具有相应保证的近似器或系数方案，再分析其谱条件。本文的条件定理不把这种实现保证作为已经证明的事实。

新增实验会在固定检查点记录 $\delta$ 的观测值，但稀疏检查点不能证明式\eqref{eq:delta-assumption}对全程成立。实际 ResNet-18 使用固定 $K$、Nesterov、形状缩放和权重衰减；其结果是实用配置的经验比较。精确版本、条件近似版本和实用混合版本的假设边界分别保留，不将一个版本的结论替换成另一个版本的保证。

\FloatBarrier
